\documentclass[10pt,a4paper]{article}

\usepackage[T1]{fontenc}
\usepackage[utf8]{inputenc}
\usepackage[ngerman]{babel}
\usepackage[a4paper,margin=14mm]{geometry}
\usepackage{lmodern}
\usepackage[table]{xcolor}
\usepackage{tabularx}
\usepackage{array}
\usepackage{microtype}
\usepackage[most]{tcolorbox}

\pagestyle{empty}
\setlength{\parindent}{0pt}
\setlength{\parskip}{0pt}
\renewcommand{\familydefault}{\sfdefault}
\renewcommand{\arraystretch}{1.12}

\definecolor{Navy}{HTML}{17324D}
\definecolor{Blue}{HTML}{3B6EA5}
\definecolor{Sky}{HTML}{EAF3FA}
\definecolor{Mint}{HTML}{E8F5EF}
\definecolor{Gold}{HTML}{F3B33D}
\definecolor{Ink}{HTML}{1E2935}
\definecolor{Muted}{HTML}{617080}
\definecolor{Line}{HTML}{D8E1E8}
\definecolor{CodeBg}{HTML}{F5F7F9}

\color{Ink}
\newcommand{\sectiontitle}[1]{%
  \vspace{2.0mm}{\color{Blue}\bfseries\large #1}\par
  \vspace{0.7mm}{\color{Blue}\rule{\linewidth}{0.7pt}}\par\vspace{1.0mm}}
\newcommand{\codebox}[1]{%
  \colorbox{CodeBg}{\parbox{\dimexpr\linewidth-2\fboxsep\relax}{\ttfamily\small #1}}}
\newcommand{\sourcecaption}[1]{\textcolor{Blue}{\bfseries\texttt{#1}}\par\vspace{0.5mm}}

\begin{document}

\colorbox{Navy}{\parbox{\dimexpr\linewidth-2\fboxsep\relax}{%
  \vspace{3mm}\color{white}\begin{tabularx}{\linewidth}{@{}Xr@{}}
  {\small\bfseries INFORMATIK · KLASSE 10} & \colorbox{Gold}{\color{Navy}\small\bfseries\strut LÖSUNGEN}\end{tabularx}
  \vspace{2.4mm}{\fontsize{18}{21}\selectfont\bfseries Aufgabe 6 -- Das Kreuzprodukt}\par
  \vspace{0.8mm}{\large Sicherungsblatt}\vspace{2.8mm}}}

\vspace{2.5mm}
\colorbox{Sky}{\parbox{\dimexpr\linewidth-2\fboxsep\relax}{\vspace{1.6mm}
\textcolor{Blue}{\bfseries Ziel:} Du bildest ein Kreuzprodukt, erklärst die Anzahl seiner Ergebniszeilen und schränkst es mit \texttt{WHERE} ein.\vspace{1.6mm}}}

\sectiontitle{1 · Vom Mensa-Angebot zur Ergebnisrelation}
{\small\color{Muted}Jede Zeile der Ergebnisrelation enthält genau eine Vorspeise, eine Hauptspeise und eine Nachspeise.}\par\vspace{1.1mm}

{\centering
\begin{tabular}{@{}c@{\hspace{7mm}}c@{\hspace{7mm}}c@{}}
\begin{minipage}[t]{48mm}\centering\sourcecaption{Vorspeise}
\begin{tabular}{|>{\ttfamily\footnotesize}l|>{\ttfamily\footnotesize\raggedleft\arraybackslash}p{13mm}|}\hline
\rowcolor{Sky}\textbf{name} & \textbf{preis}\\\hline
Lauchsuppe & 1,50~€\\\hline
Salat & 2,00~€\\\hline
Tagessuppe & 1,00~€\\\hline
Rohkost & 1,35~€\\\hline
\end{tabular}\end{minipage}
&
\begin{minipage}[t]{48mm}\centering\sourcecaption{Hauptspeise}
\begin{tabular}{|>{\ttfamily\footnotesize}l|>{\ttfamily\footnotesize\raggedleft\arraybackslash}p{13mm}|}\hline
\rowcolor{Sky}\textbf{name} & \textbf{preis}\\\hline
Käsespätzle & 3,50~€\\\hline
Reispfanne & 2,50~€\\\hline
Pizza & 3,44~€\\\hline
\end{tabular}\end{minipage}
&
\begin{minipage}[t]{48mm}\centering\sourcecaption{Nachspeise}
\begin{tabular}{|>{\ttfamily\footnotesize}l|>{\ttfamily\footnotesize\raggedleft\arraybackslash}p{13mm}|}\hline
\rowcolor{Sky}\textbf{name} & \textbf{preis}\\\hline
Gemischtes Eis & 2,50~€\\\hline
\end{tabular}\end{minipage}
\end{tabular}\par}

\vspace{0.5mm}
{\centering\textcolor{Blue}{\bfseries Ausgangstabellen \; $\longrightarrow$ \; Kreuzprodukt \; $\longrightarrow$ \; Ergebnisrelation}\par
\vspace{0.5mm}\colorbox{Mint}{\parbox{74mm}{\centering\textbf{4 · 3 · 1 = 12 Datensätze}}}\par}
\vspace{0.8mm}

{\small\textcolor{Blue}{\bfseries Ergebnisrelation: alle 12 Menüs}}\par\vspace{0.4mm}
{\footnotesize\ttfamily
\begin{tabularx}{\linewidth}{@{}>{\raggedright\arraybackslash}X>{\raggedleft\arraybackslash}p{15mm}>{\raggedright\arraybackslash}X>{\raggedleft\arraybackslash}p{15mm}>{\raggedright\arraybackslash}X>{\raggedleft\arraybackslash}p{15mm}@{}}
\rowcolor{Sky}\multicolumn{2}{c}{\bfseries Vorspeise} & \multicolumn{2}{c}{\bfseries Hauptspeise} & \multicolumn{2}{c}{\bfseries Nachspeise}\\
\rowcolor{CodeBg}\textbf{name} & \textbf{preis} & \textbf{name} & \textbf{preis} & \textbf{name} & \textbf{preis}\\
Lauchsuppe & 1,50~€ & Käsespätzle & 3,50~€ & Gemischtes Eis & 2,50~€\\
Lauchsuppe & 1,50~€ & Reispfanne & 2,50~€ & Gemischtes Eis & 2,50~€\\
Lauchsuppe & 1,50~€ & Pizza & 3,44~€ & Gemischtes Eis & 2,50~€\\
Salat & 2,00~€ & Käsespätzle & 3,50~€ & Gemischtes Eis & 2,50~€\\
Salat & 2,00~€ & Reispfanne & 2,50~€ & Gemischtes Eis & 2,50~€\\
Salat & 2,00~€ & Pizza & 3,44~€ & Gemischtes Eis & 2,50~€\\
Tagessuppe & 1,00~€ & Käsespätzle & 3,50~€ & Gemischtes Eis & 2,50~€\\
Tagessuppe & 1,00~€ & Reispfanne & 2,50~€ & Gemischtes Eis & 2,50~€\\
Tagessuppe & 1,00~€ & Pizza & 3,44~€ & Gemischtes Eis & 2,50~€\\
Rohkost & 1,35~€ & Käsespätzle & 3,50~€ & Gemischtes Eis & 2,50~€\\
Rohkost & 1,35~€ & Reispfanne & 2,50~€ & Gemischtes Eis & 2,50~€\\
Rohkost & 1,35~€ & Pizza & 3,44~€ & Gemischtes Eis & 2,50~€\\
\end{tabularx}}

\sectiontitle{2 · Kreuzprodukt mit SQL}
\begin{minipage}[t]{0.37\linewidth}
\vspace{0pt}
\textcolor{Blue}{\bfseries Vollständiges Kreuzprodukt}\par\vspace{0.8mm}
\codebox{SELECT *\par FROM Vorspeise, Hauptspeise, Nachspeise;}
\end{minipage}\hfill
\begin{minipage}[t]{0.60\linewidth}
\vspace{0pt}
\begin{tcolorbox}[colback=Mint,colframe=Blue,arc=2mm,boxrule=.8pt,left=2.5mm,right=2.5mm,top=1.5mm,bottom=1.5mm]
\textbf{Merksatz}\par\smallskip
Wenn bei einer \texttt{SELECT}-Abfrage beim Schlüsselwort \texttt{FROM} mehrere Tabellen angegeben werden, wird das Kreuzprodukt auf diesen Tabellen angewendet.\par\smallskip
Beim Kreuzprodukt wird jeder Datensatz der einen Tabelle mit jedem Datensatz der anderen Tabelle (bzw. den Datensätzen aller weiteren Tabellen) verknüpft.
\end{tcolorbox}
\end{minipage}

\sectiontitle{3 · Mit \texttt{WHERE} einschränken}
\begin{minipage}[t]{0.48\linewidth}
\textcolor{Blue}{\bfseries Salat als Vorspeise}\par\vspace{0.7mm}
\codebox{SELECT *\par FROM Vorspeise, Hauptspeise, Nachspeise\par WHERE Vorspeise.name = 'Salat';}\par\vspace{0.7mm}
\textbf{Ergebnis: 3 Ergebniszeilen}
\end{minipage}\hfill
\begin{minipage}[t]{0.48\linewidth}
\textcolor{Blue}{\bfseries Pizza als Hauptspeise}\par\vspace{0.7mm}
\codebox{SELECT *\par FROM Vorspeise, Hauptspeise, Nachspeise\par WHERE Hauptspeise.name = 'Pizza';}\par\vspace{0.7mm}
\textbf{Ergebnis: 4 Ergebniszeilen}
\end{minipage}

\vspace{1.2mm}
{\small\textcolor{Muted}{\textbf{Warum der Tabellenname?} \texttt{name} kommt in mehreren Tabellen vor. Die Angabe \texttt{Vorspeise.name} beziehungsweise \texttt{Hauptspeise.name} macht die Bedingung eindeutig.}}\par\vspace{1.0mm}
\colorbox{Sky}{\parbox{\dimexpr\linewidth-2\fboxsep\relax}{\vspace{1.3mm}\small
\textcolor{Blue}{\bfseries Filterbedingung:} Die Bedingung \texttt{Vorspeise.name = 'Salat'} filtert alle Datensätze heraus, in denen eine Vorspeise vorkommt, die nicht Salat entspricht.\vspace{1.3mm}}}

\vfill{\color{Line}\rule{\linewidth}{.5pt}}\par\vspace{1.5mm}
\begin{tabularx}{\linewidth}{@{}Xr@{}}{\small\color{Muted}Klasse 10 · Datenbanken} & {\small\color{Muted}Sicherungsblatt · Aufgabe 6}\end{tabularx}
\end{document}
